\appendix

\chapter{Függelékek}
\label{chap:fuggelek}

A függelék tartalmazza az „Automatikus hangszerfelismerés zenefelvételek alapján” című államvizsga-szakdolgozat teljes forráskódját, az adathalmaz-struktúrát és a betanított modellfájlokat tartalmazó digitális adathordozó könyvtárszerkezetének leírását, valamint a felhasznált szoftverfüggőségek összegzését.

%% ═══════════════════════════════════════════════════════════════════════════
\section{A projekt könyvtárszerkezete}
\label{sec:konyvtarszerkezet}

A benyújtott digitális függelék (amely a GitHubon elérhető nyilvános Git adattárban\footnote{\url{https://github.com/nCsab/Instrument-Recognition}} is megtalálható) az alábbi főbb könyvtárakat tartalmazza. A számozott szkriptek a feldolgozási csővezeték (pipeline) egymásra épülő lépéseit jelzik.

\begin{table}[ht!]
\centering
\small
\caption{A projekt legfelső szintű könyvtárszerkezete}
\label{tab:konyvtarszerkezet}
\begin{tabular}{p{4.5cm}p{10cm}}
\toprule
\textbf{Könyvtár / Fájl} & \textbf{Leírás} \\
\midrule
\texttt{thesis/} & A \LaTeX\ formátumú szakdolgozat forrásfájljai, stílusfájl (\texttt{sapientia-thesis.cls}), irodalomjegyzék (\texttt{references.bib}) és a generált \texttt{main.pdf}. \\
\texttt{scripts/} & A hangszerfelismerő rendszer Python szkriptjei (részletezés alább). \\
\texttt{owndataset/} & A manuálisan összeválogatott és kivágott internetes hangminták hangszerosztályonként külön almappákban, valamint az újrafelvételezett minták a \texttt{recorded\_from\_mic/} almappában. \\
\texttt{hybrid\_dataset\_\allowbreak own\_final/} & A végleges, összeszerelt hibrid adathalmaz 7~osztálymappával (\texttt{guitar}, \texttt{piano}, \texttt{vocal}, \texttt{string}, \texttt{reed}, \texttt{brass}, \texttt{noise}), osztályonként 1500~db, 1~másodperces, 16\,kHz-es WAV minta. \\
\texttt{processed\_data/} & Az előállított jellemzőmátrixok NumPy formátumban: \texttt{X\_stft\_full.npy}, \texttt{X\_log\_mel\_full.npy}, \texttt{X\_mfcc\_full.npy} és a közös címkefájl \texttt{y\_labels\_full.npy}. \\
\texttt{models/} & A betanított Keras modellfájlok: \texttt{best\_stft\_2dcnn\_model.keras}, \texttt{best\_log\_mel\_2dcnn\_model.keras}, \texttt{best\_mfcc\_2dcnn\_model.keras}. \\
\texttt{augmented\_samples/} & Az augmentációs lánc bemutató mintái: minden hangszerosztályból 3-3 eredeti--augmentált páros WAV fájl ellenőrzés céljából. \\
\texttt{dataset/} & A letöltött nyilvános referencia-adathalmazok (IRMAS, NSynth, Medley-solos-DB, TinySOL, Good-sounds, ESC-50) az elemzéshez. \\
\bottomrule
\end{tabular}
\end{table}
\clearpage
%% ═══════════════════════════════════════════════════════════════════════════
\section{A feldolgozási csővezeték szkriptjei}
\label{sec:pipeline_scripts}

A rendszer teljes feldolgozási lánca indexelt Python szkriptekből áll, amelyek a tervezés sorrendjében futtathatók. Az alábbi táblázat leírja minden szkript feladatát.

\begin{table}[ht!]
\centering
\small
\caption{A feldolgozási csővezeték szkriptjei és feladataik}
\label{tab:pipeline_scripts}
\begin{tabular}{p{5.8cm}p{8.7cm}}
\toprule
\textbf{Szkript} & \textbf{Feladat} \\
\midrule
\texttt{01\_analyze\_datasets.py} & A nyilvános referencia-adathalmazok (IRMAS, NSynth stb.) szerkezetének és statisztikáinak feltárása, összehasonlító elemzés készítése. \\
\texttt{02\_mic\_prep\_\allowbreak concatenate.py} & Az újrafelvételezéshez szükséges hanganyag előkészítése: a kiválasztott minták összefűzése egyetlen hosszú WAV fájlba sípszó-szinkronizálással. \\
\texttt{03\_mic\_record\_audio.py} & A laptop beépített mikrofonjával vagy mobiltelefon mikrofonjával történő visszajátszás és valós idejű rögzítés. \\
\texttt{04\_mic\_slice\_audio.py} & A rögzített mikrofonos felvétel darabolása 1~másodperces szegmensekre a sípszó-szinkronizáció alapján. \\
\texttt{05\_assemble\_dataset.py} & A végleges hibrid adathalmaz összeállítása: a tiszta minták, az újrafelvételek és az augmentált minták egyesítése és randomizált másolása a \texttt{hybrid\_dataset\_own\_final/} könyvtárba. \\
\texttt{05a\_prepare\_hybrid\_\allowbreak generic.py} & Kiegészítő előkészítő szkript a hibrid adathalmaz generikus felépítéséhez. \\
\texttt{06\_extract\_features.py} & A fő jellemzőkinyerő szkript, amely az összes hangmintából STFT, log-mel és MFCC spektrogramokat generál, és NumPy mátrixokba menti. \\
\texttt{06\_extract\_features/} & Alkönyvtár a három jellemzőtípus önálló kinyerő szkriptjeivel: \texttt{06\_extract\_stft.py}, \texttt{06\_extract\_log\_mel.py}, \texttt{06\_extract\_mfcc.py}. \\
\texttt{07\_realtime\_\allowbreak recognition.py} & A valós idejű hangszerfelismerő modul: mikrofon-bemenet, csúszóablak, előrejelzés és megjelenítés. \\
\texttt{10\_train\_model\_colab/} & A Google Colab felhőkörnyezetben futtatott tanító szkriptek: \texttt{10\_train\_stft.py}, \texttt{10\_train\_log\_mel.py}, \texttt{10\_train\_mfcc.py}. \\
\texttt{utils/} & Segédmodulok: \texttt{augmentation\_utils.py} (az augmentációs lánc függvényei) és \texttt{feature\_utils.py} (jellemzőkinyerési segédfüggvények). \\
\bottomrule
\end{tabular}
\end{table}

\clearpage
%% ═══════════════════════════════════════════════════════════════════════════
\section{Fejlesztői és futtatási környezet}
\label{sec:fejlesztoi_kornyezet}

Az alábbi táblázat tartalmazza a projekt fejlesztéséhez és futtatásához szükséges szoftverfüggőségeket.

\begin{table}[H]
\centering
\caption{A fejlesztői környezet főbb összetevői}
\label{tab:fejlesztoi_kornyezet}
\begin{tabular}{p{4.2cm}p{2.5cm}p{7.5cm}}
\toprule
\textbf{Szoftver / Könyvtár} & \textbf{Verzió} & \textbf{Felhasználás} \\
\midrule
Python & 3.10+ & A teljes feldolgozási csővezeték nyelve \\
\addlinespace
TensorFlow / Keras & 2.x & A 2D CNN modell építése, tanítása és mentése \\
\addlinespace
Librosa & 0.10+ & Hangfájlok betöltése, STFT, mel-spektrogram és MFCC kinyerés \\
\addlinespace
NumPy & 1.24+ & Jellemzőmátrixok tárolása és manipulálása \\
\addlinespace
SoundFile / SoundDevice & -- & Hangfájlok olvasása/írása és valós idejű mikrofon-hozzáférés \\
\addlinespace
Matplotlib & 3.x & Diagramok, spektrogramok és kiértékelési ábrák készítése \\
\addlinespace
Google Colab & -- & GPU-gyorsított modelltanítás felhőkörnyezetben \\
\addlinespace
Git & 2.x & Verziókezelés és a projekt előrehaladásának dokumentálása \\
\addlinespace
\LaTeX\ (TeX Live 2026) & -- & A szakdolgozat szedése a Sapientia sablon alapján \\
\bottomrule
\end{tabular}
\end{table}
